\documentclass[11pt,a4paper]{article}

% Engineering debug report. An equal rank deliverable: the record of the fights
% that the clean results hide. No en or em dashes; enforced by the dash check.

\usepackage[utf8]{inputenc}
\usepackage[T1]{fontenc}
\usepackage{lmodern}
\usepackage[margin=1in]{geometry}
\usepackage{amsmath}
\usepackage[hidelinks]{hyperref}
\usepackage{xcolor}

\newcommand{\entry}[5]{%
  \subsection*{#1}%
  \noindent\textbf{Symptom.} #2\par\smallskip
  \noindent\textbf{Root cause.} #3\par\smallskip
  \noindent\textbf{Fix and why.} #4\par\smallskip
  \noindent\textbf{Verification.} #5\par\bigskip
}

\title{Engineering Debug Report\\
\large CPU Microbenchmark Suite}
\author{Olajide Badejo}
\date{July 2026}

\begin{document}
\maketitle

\section{Purpose}

This report records the debugging that the clean results in the main report do
not show. Each entry states the symptom, the root cause, the fix and why it was
chosen over alternatives, and how the fix was verified. The anti optimization
fight belongs here, and on this platform it took the form of a clock measurement
fight rather than a deleted loop. This document is an equal rank deliverable.

\section{Findings}

\entry{The topology hides the hybrid P and E split}
{The topology reader could not tell performance cores from efficiency cores on
the target. The sysfs paths that expose the hybrid split do not exist under
WSL2, and every logical CPU reports the same 2 MB L2, so the split cannot be
recovered from cache size either.}
{WSL2 runs under a Hyper-V utility virtual machine that flattens the CPU
topology it presents to the guest. The hybrid enumeration is a bare metal kernel
feature the virtualized sysfs does not carry through.}
{The reader detects the split when the sysfs files are present and marks it as
known; when absent it falls back to the documented Intel enumeration and marks
the result as heuristic, so no caller mistakes the guess for a measurement. The
real parsing path is still tested against a committed hybrid fixture so the bare
metal code cannot rot.}
{The topology test passes both a flat WSL style fixture and a hybrid bare metal
style fixture, asserting 16 performance and 12 efficiency cores in the hybrid
case.}

\entry{The reported clock does not track turbo}
{The first pointer chase runs reported L1 latency at about 3.2 cycles, which is
impossible: the Raptor Cove L1 latency is 5 cycles. The nanosecond timings were
right, but the cycle conversion was too low.}
{Cycles were computed from the \texttt{cpu MHz} field in
\texttt{/proc/cpuinfo}, which under WSL2 is pinned at about 3.42 GHz and never
moves, even under load, while the pinned core actually boosts to about 5.3 GHz.
There is no cpufreq interface to consult either.}
{A software clock measurement was added: a two billion iteration dependent
integer add chain runs at one cycle per iteration, so frequency is iterations
over elapsed time. It needs no counter and follows turbo. The wrong cpuinfo
value is still recorded for transparency.}
{The calibration reports about 5.3 GHz under load, and with it L1 reads 5.0
cycles, L2 16.3, and DRAM 85 to 93 ns, all within the sanity gates. A later
single shot instability was fixed by warming up and taking the maximum over four
passes, after a lone 3.81 GHz reading produced an impossible 138 percent FLOPS
efficiency.}

\entry{The plateau detector mixed log and raw error}
{On a synthetic four plateau curve the change point detector returned the wrong
boundaries, splitting the flat DRAM plateau instead of the obvious L1 to L2
step.}
{Change point gains were computed in log latency space, the right space, but the
child segments after a split had their error recomputed in raw space. A DRAM
child then carried a raw error near 65 while every candidate gain was a log space
quantity near 2, so a spurious gain of about 65 beat the real L1 to L2 split.}
{Compute the child segment error in the same log space as the gains. One line.}
{The synthetic test recovers the true change points within one index and the
plateau means within 25 percent. On the real curve the L1 and L2 boundaries land
within one octave of the sysfs capacities.}

\entry{A main thread barrier race gave impossible bandwidth}
{The STREAM scaling sweep reported thousands of GB/s, physically impossible on a
dual channel desktop, and disagreed with the single thread path even at one
thread.}
{The parallel region was timed on the main thread across two barriers, and a
subtle phase interaction let the timer capture a window in which little work had
run, so the minimum over trials latched onto a near zero time.}
{Move all timing into the workers: each worker times its own chunk between two
barriers, and worker zero reduces to the slowest worker per trial and keeps the
minimum. The barriers provide the happens before needed to read every worker's
time safely, and there is no main thread window to race.}
{Bandwidth became physically sane: saturating near 78 GB/s, 87 percent of the
theoretical dual channel peak, and the single thread path and the one thread
scaling point now agree.}

\entry{Unsigned underflow poisoned the GEMM test data}
{The GEMM correctness test failed on every element, yet the blocked result
matched the reference exactly. The values were infinity.}
{The data was initialized with \texttt{(i \% 11) - 5} where \texttt{i} is
unsigned, so residues below 5 underflowed to about 1.8e19 and the products
overflowed to infinity. Both kernels agreed at infinity, but the near check
computes the absolute difference of two infinities, which is not a number, so
every comparison failed.}
{Cast the residue to a signed integer before subtracting. One line in each of
the test and the validator.}
{The GEMM test passes for five tile sizes including tiles that do not divide the
matrix dimensions.}

\entry{The BLIS prediction misses the naive kernel, by design}
{The BLIS analytical model predicts a large block $K_c$ near 700, but the naive
blocked kernel's best tile uses $K_c$ near 64, a gap of about 3.4 octaves.}
{The model assumes a packed micro kernel whose L1 resident panel of $B$ is only
$N_r$ elements wide. The naive kernel does not pack and its inner loop spans the
full width of $B$, so each reused row is the full matrix width and only a few
fit in cache, pushing the optimum to a much smaller $K_c$.}
{No code fix: this is the reported finding for objective 4. The main report
presents both tiles and attributes the gap to the absence of packing. A packed
micro kernel is recorded as future work.}
{The tile prediction script prints the predicted and empirical tiles and the
octave gap, labeled a miss. The roofline shows the same gap from the performance
side: the naive GEMM sits below the single core ceiling.}

\entry{No QEMU floating point quirk, and why that is worth stating}
{The NEON port was expected, from folklore, to show a floating point
discrepancy under qemu-user.}
{None appeared. NEON \texttt{vfmaq\_f32} is a true fused multiply add with a
single rounding, matching the x86 FMA, so the same numerical tolerances hold on
both.}
{The width independent FLOPS correctness test and all three NEON STREAM kernel
variants pass under qemu-user with the same tolerances as the native build. The
absence of a quirk is recorded so the next reader does not go looking for one.}

\end{document}
